% =====================================================================
% overleaf/03_architecture.tex — LaTeX rendering of en/03_architecture.md v0.7 (17 September 2026)
% Inline math only (amsmath is loaded by aamas.cls): the agent tuple, the policy vector and the
%   categorical draw. No display equation, no numbered equation. The action space is deliberately
%   kept in plain words.
% Dernière mise à jour : 2026-09-17
% Chapter 3 of the paper, a \section of its own, to be inserted right after overleaf/02_related_work.tex.
% This chapter cites nothing: no \citep / \citet, and it adds no entry to references.bib.
% Cross-references use \label / \ref, so no number is hard-coded.
% Forward reference: sec:results (chapter 6) is NOT defined yet, that chapter having no LaTeX
%   rendering so far. It resolves to ?? until 06 is written; the label is reserved with this name.
%   sec:metrics is defined by overleaf/04_metrics_and_substrate.tex: keep the label spelling exactly.
% Ready to paste into the Overleaf project (acmart + natbib: \citep, \citet).
% Derived file: regenerate after every validated pass of en/03_architecture.md.
% Preamble additions:
%   \newcommand{\todo}[1]{\textbf{[#1]}}   % placeholder figures, filled from methode/experience_plan/experiments.yaml
% FIGURE — upload docs/paper/article/images/architecture_GAMA_Agents.jpg and ch3_oubli.png into the
%   Overleaf project under images/. graphicx is already loaded by aamas.cls (l. 593): nothing to add
%   to the preamble. ch3_oubli.png is regenerated by scripts/analysis/plot_chapitre3.py, which reads
%   the constants from the code; re-upload it after any change to the memory settings.
%   The template requires the caption BELOW the figure and a mandatory \Description{} (plain text,
%   2000 characters at most, not printed). Set as figure* (both columns): the diagram is 16:9, and a
%   single column would make it unreadable. Switch to figure + \columnwidth if the page budget demands it.
% Comments "% source:" give the repository file a plain-text figure is recomputed from.
% =====================================================================
% =====================================================================


\section{Architecture of the agentic device}
\label{sec:system}

\subsection{Overview}
\label{sec:overview}

The device couples three components. A GAMA multi-agent simulation carries the world: the real geography of the Toulouse area, the networks, the clock, and the physical execution of trips. A controller carries the agent life cycle: it detects that an agent is due to leave and builds the itinerary supply actually available at that hour. A decision module carries the choice itself: it is there, and only there, that the language model intervenes.

\begin{figure*}[t]
  \centering
  \includegraphics[width=\textwidth]{architecture_GAMA_Agents.jpg}
  \caption{The device. At the bottom, the GAMA simulation carries the bodies, the network and the clock; at the top, the generative agent carries memory, reasoning and decision; between them, the exchange pipeline: perception goes up over HTTP, actions come back down over WebSocket. On the right, the two engines that produce the itinerary supply; on the left, the gateway that serves the calls to the language model.}
  \Description{Block diagram of the device, on three horizontal bands. Top band, the generative agent: a group of agent icons connected on the left to an LLM gateway box (provider rotation, cache), on the right to two boxes, daily activity planning and travel planning, and below to two memory stores, short-term memory and long-term memory, linked by an arrow labelled reflection; a double arrow between the agents and the memories is labelled memory read and write. Two dashed boxes feed the agent from above: population data (personal identity, daily activities) and environment (weather, traffic). Middle band, a horizontal bar labelled data exchange pipeline, with Python services above and GAMA in Java below; perception data goes up over HTTP to the short-term memory, actions come down over WebSocket. Bottom band, the GAMA platform simulation: a map of the Toulouse network and icons for bus, train, tram, cable car, bicycle, pedestrian and car, fed from the right by GTFS and OSM map data. On the right of the figure, a dashed box named routing engines contains OpenTripPlanner for public transport and OSMnx for walking, cycling and car, with results cached; it receives requests from travel planning and returns itineraries. At the foot, two summary boxes: simulation, Agent GAMA, physical body, position, network; and cognition, Python server with the language model, memory, reasoning, decision; between them the two channels, HTTP for state and WebSocket for action and perception.}
  \label{fig:architecture}
\end{figure*}

The loop is closed. An idle agent triggers a planning step; it receives the itineraries that exist for that trip and spreads its preference across them; its decision goes to the simulation, which executes the trip on the network and returns what actually happened: transfers, waits, arrival time. These returns feed the agent's memory, which will weigh on the next decision.

A thousand agents decide asynchronously while the simulated clock runs. The controller therefore serves planning tasks by increasing deadline, the deadline being the simulated departure time, and holds the GAMA clock whenever the estimated time to drain the queue threatens a deadline. No agent leaves before it has decided, and no decision is taken after the fact.

Formally, an agent is a quadruple $\text{Agent}_i = \langle P_i, M_{i,t}, C_{i,t}, \pi_\theta \rangle$: $P_i$ the sealed sociodemographic persona, restricted to the variable contract of Section~\ref{sec:decision}; $M_{i,t}$ the two-component memory register at time $t$, a circular buffer of recent observations and a long-term vector index (Section~\ref{sec:memory}); $C_{i,t}$ the state at time $t$ of the household vehicle chain (where the car and the bicycle are, who holds a licence, Section~\ref{sec:chain}); and $\pi_\theta$ the decision policy verbalised by the language model $\theta$. The first three terms tell agents apart; the fourth is common to them all, and it is the one the rest of the paper puts to the test.

The agent has a fixed schedule: the same chain of activities, in the same order, every weekday, closing in the evening on the return home. Weekends carry no activity, for simplicity. It must carry that schedule out with what the city offers at that hour, in that place, and with the vehicles it has to hand. Those bounds are set by the controller, before the question is ever put to the model. What the simulation adds is of another nature: it produces the consequences. The trip as lived is not the trip as planned; a missed bus, a missed transfer, a jam born of an incident, a quarter of an hour of delay are facts produced by execution, and it is these that come back into the loop. This is what separates the device from a classifier queried in a vacuum, and what makes the floor and the ceiling of Section~\ref{sec:metrics} comparable with the agent: all three decide on the same substrate.

\subsection{The itineraries offered}
\label{sec:field}

The itineraries offered to the agent come from two sources. OpenTripPlanner produces the public transport itineraries: bus, tram, metro, regional trains, intercity coaches, on the real timetables of the networks. A shortest-path computation on the OpenStreetMap graph produces the direct trips: walking, cycling, car. The perimeter is that of the reference survey, the Cerema's EMC\textsuperscript{2} 2023: the 453 municipalities of the Toulouse metropolitan area.

Not every mode is offered to everyone, and that restriction operates before the decision. Car and bicycle are only proposed if the agent owns them and they are at the point of departure. Driving further requires a licence and legal age, but an agent who does not drive may still travel by car as a \textbf{passenger} in a motorised household, and the option stays open on that footing.
% source: docs/arch/routing.md, § Modes disponibles par agent ; docs/arch/vehicle-chain.md, § 1 bis

\subsection{The decision point}
\label{sec:decision}

At the moment of choosing, the agent receives a short profile (first name, age, occupation, household composition, income level), the destination and its zone, the departure time, the current weather and that of the remaining slots of the day, the trips it still has to make before returning home, the memories judged relevant to this trip, and finally the numbered list of itineraries, each with its mode and the detail of its legs. What it can do is not an abstract set of modes: it is the finite set of itineraries viable at that instant, $\mathcal{A}(o_t, C_{i,t}) \subseteq \mathcal{O}(o_t)$, where $\mathcal{O}(o_t)$ is what the two engines of Section~\ref{sec:field} produce for the observation $o_t$, and where the inclusion becomes strict as soon as eligibility or the position of the vehicles $C_{i,t}$ removes an option (Section~\ref{sec:chain}).

The language model does not designate an itinerary: it \textbf{spreads a probability mass over the options}, one entry per option, each with its justification. This is the policy $\pi_\theta$ over that list: it yields a vector $\mathbf{p}_t = \pi_\theta(\cdot \mid o_t, M_{i,t})$ over the simplex of viable options, where $o_t$ is the observation of the trip. The decision is a categorical draw from that vector, $a_t \sim \mathrm{Cat}(\mathbf{p}_t)$, seeded from that same context, reproducible identically, and never an $\arg\max$: this is what separates the device from a classifier, and what leaves it its individual heterogeneity. The options are presented to it in a randomised order so that their rank does not become a preference. This detail settles the nature of what we measure: what the model produces is a distribution.

\subsection{Memory}
\label{sec:memory}

Without memory, an agent decides identically every morning, and no dynamics can be observed. The device carries two levels of it.

Short-term memory is a circular buffer per agent, a hundred entries at most, partitioned by activity. It does not record decisions alone: it receives the physical experience the simulation produces: transfers, waits at the stop, arrival times, the events of the day. This is how lived experience enters the reasoning, and not through some description the agent would give of itself.

Consolidation fires on two conditions: as soon as the buffer reaches ten entries, and in any case once per simulated day, late in the evening.

It writes into long-term memory, a vector index specific to each agent, where two registers coexist without obeying the same clock. \textbf{Episodic} traces (what happened, and when) erode: their weight decays exponentially from the last recall, with a time constant of 2.8 days that a serious memory lengthens, never beyond thirty. \textbf{Concepts} (what the agent holds to be true, ``this line saturates when it rains'') do not erode: they live on a confidence that each observation confirms or contradicts, a fact not becoming false because ten days have passed. This is Tulving's distinction between episodic and semantic memory, and it is what settles what gets forgotten (Figure~\ref{fig:memory-clocks}).

\begin{figure*}[t]
  \centering
  \includegraphics[width=\textwidth]{ch3_oubli.png}
  \caption{The temporal component of the retrieval score, by age since last recall. Episodic traces at two levels of gravity; a concept before and after five counter-examples.}
  \Description{Line chart. The horizontal axis is age since last recall, zero to thirty days; the vertical axis is the temporal component of the retrieval score, zero to one. Two descending curves start at one: an ordinary episodic trace with a lifetime of 2.8 days, which reaches 0.5 at 1.9 days and is purged at 13 days, and a serious episodic trace with a lifetime of 19.6 days, still at 0.22 on day thirty; the area between them is shaded. A dotted curve between the two is the sensitivity arm at 8.3 days. Two horizontal lines carry the concepts: 0.93 for a concept with twelve observations and no counter-example, 0.68 for the same concept after five counter-examples, with an arrow between them. A dash-dotted line at 0.5, in the colour of the concepts and spanning only their part of the chart, marks where a concept is set aside, having been contradicted more often than confirmed.}
  \label{fig:memory-clocks}
\end{figure*}

Retrieval is not a plain similarity search: it combines five terms: semantic proximity, the gravity of the memory, time decay, the match on weather, and affinity of place, time slot and purpose over an age window equal to the horizon of the experiment, sixty days at most. And the model does not receive raw memories. It receives what the agent has aggregated from them: what it does most often, what it holds to be true, what has just changed. Three memories chosen for this trip are added to that.


The agent does not merely remember that it arrived late: it draws a consequence from it. A late arrival brings forward its target time for that activity, by 75\,\% of the observed delay, never encroaching on the preceding or the following activity. The readjustment is written into its memory in the first person, so that it can read it again when deciding tomorrow. This is the most elementary form of learning the device carries, and it owes nothing to the language model: it is produced by the gap between the planned schedule and the schedule as lived.

\subsection{Trip-chain modelling}
\label{sec:chain}

An agent who cycled to work has no car at the office in the evening. That obvious fact is not obvious to a system that would decide each trip independently, and it is what separates the device from a classifier called once per trip.

The controller therefore holds, for each agent, the position of their personal vehicles, and draws three rules from it: a vehicle mode is only proposed if the vehicle is parked at the point of departure; the vehicle of the chosen mode follows its user to the destination, the others stay where they are; and on a return trip, if a vehicle is waiting at the origin, the candidate itineraries are restricted to that mode. This last rule filters the options, it does not add a decision.
% source: docs/arch/vehicle-chain.md, § Les trois règles

The methodological consequence has to be drawn, and Section~\ref{sec:results} returns to it: part of the simulated modal split is produced \textbf{under} these constraints, so as to be more representative of real situations. An agent whose car stayed at home does not choose not to drive, it cannot. This is why the floor and the ceiling of Section~\ref{sec:metrics} are placed under exactly the same constraints: without that, we would be comparing decisions taken in different worlds.

% =====================================================================
% TICKETS ATTACHED TO THIS CHAPTER — not part of the text, kept in step with en/ and fr/03_architecture.md
%   Ticket 050 — mathematical formalisation of the agentic device: APPLIED in v0.5 for the tuple
%     ("Overview") and the draw ("The decision point"). Its second criterion, the action space, is
%     given in plain words rather than in notation: the author's decision of 2026-09-15.
%   Ticket 051 — reflection and overhaul of the memory architecture.
%   Ticket 070 — accidents drawn at random on the network and the delay they inflict: the mechanism
%     promised by "a jam born of an incident" in "Overview". Written ahead of the code, by intent.
%   Ticket 071 — evolution of memory, from the current code to the target state: "Memory" describes
%     that target state, likewise ahead of the code.
%
% OPEN DEPENDENCY
%   Ticket 046 — informational asymmetry of the evaluation contract. "The decision point" says that
%     the agent and the statistical models see the same thing; the ticket establishes that the agent
%     receives more. To be revisited once the ticket is settled. Noted, not corrected.
% =====================================================================
